The simplest Bell experiment allows two binary measurements on each of two separated systems. We prove that its complete quantum correlation set, including all marginal expectations and joint correlators, admits a finite semidefinite representation. Consequently, one fixed semidefinite program suffices to optimize every Bell functional in this scenario exactly. The proof reduces arbitrary-dimensional quantum realizations to real two-qubit correlations and constructs exact sums of squares through the singular limits in which measurements become compatible. The inequality coefficients and measurement parameters vary independently throughout the argument. A real-algebraic criterion converts these certificates into a finite global representation, whose matrices can be chosen with real algebraic entries.
